\chapter{Wyzwania i problemy implementacyjne}

Stworzenie aplikacji łączącej konwencjonalne generowanie danych z wykorzystaniem modeli sztucznej inteligencji wiązało się z problemami. W tym rozdziale opisano napotkane trudności techniczne oraz sposoby ich rozwiązania.

\section{Informowanie użytkownika o postępie}
Jednym z istotnych problemów na początku prac było płynne wyświetlanie paska postępu generowania na stronie internetowej, podczas gdy cały proces odbywał się w tle na serwerze.

\subsection*{Problem z HTTP}
Standardowe zapytania \gls{HTTP} nie sprawdzają się najlepiej w tym zastosowaniu, ponieważ przeglądarka musiałaby nieustannie pytać serwer o aktualny status. Przy większej liczbie użytkowników spowodowałoby to znaczne przeciążenie sieci.

\subsection*{Rozwiązanie}
Zdecydowano się na wykorzystanie protokołu WebSocket. W backendzie utworzono specjalny endpoint \texttt{/ws/jobs/\{job\_id\}}, do którego podłącza się aplikacja. Program działający w tle regularnie zapisuje swój postęp do pamięci współdzielonej, a połączenie WebSocket odczytuje te dane co pół sekundy i przesyła na bieżąco do przeglądarki. Dzięki temu interfejs działa responsywnie i płynnie aktualizuje pasek postępu.

\section{Wymuszanie unikalności danych z modelu LLM}
Modele językowe charakteryzują się niedeterministycznym działaniem, co stanowiło wyzwanie w przypadku generowania kolumn, dla których zaznaczono wymóg unikalności (np. dla numerów seryjnych generowanych przez AI).

\subsection*{Powtarzalność wyników}
Podczas testów zauważono, że przy ustawieniu wysokiej ,,temperatury'' (parametr \texttt{temperature}) model często tracił kontekst i nie trzymał się narzuconego formatu. Z kolei przy niskiej temperaturze wykazywał tendencję do generowania tych samych słów i schematów.

\subsection*{Mechanizm ponowień}
Aby zapobiec powstawaniu duplikatów, zaimplementowano własny mechanizm weryfikacji:
\begin{enumerate}
    \item Aplikacja zapisuje w pamięci listę wygenerowanych wartości dla danej unikalnej kolumny.
    \item Przy każdej nowej odpowiedzi od sztucznej inteligencji następuje sprawdzenie, czy dana wartość nie została już wykorzystana.
    \item W przypadku wykrycia duplikatu, system ponawia zapytanie (maksymalnie do 10 razy). Do promptu doklejana jest dodatkowa instrukcja: \texttt{CONSTRAINT: Value MUST be unique. DO NOT use: [lista wykorzystanych wartości]}.
    \item Aby pomóc modelowi w znalezieniu nowej odpowiedzi, przy kolejnych próbach aplikacja delikatnie zwiększa parametr \texttt{temperature} (+0.1), co zachęca algorytm do większej różnorodności w doborach słów.
\end{enumerate}

\section{Zależności cykliczne}
Aplikacja sortuje tabele przed generowaniem, aby zapewnić odpowiednią kolejność ich tworzenia (np. najpierw użytkownicy, potem ich opinie). Problem pojawiał się, gdy Tabela A wymagała klucza z Tabeli B, a Tabela B jednocześnie wymagała klucza z Tabeli A.

Tworzył się wówczas tzw. cykl zależności, który zapętlał program. Na obecnym etapie przyjęto założenie, że użytkownik musi świadomie projektować strukturę bazy i unikać takich sytuacji w interfejsie. W przyszłości planowane jest dodanie dodatkowej walidacji, która blokowałaby zapis schematu w przypadku wykrycia ,,błędnego koła''.

\section{Optymalizacja wydajności generowania przez modele lokalne}
\label{sec:llm_optimization}

Największym wyzwaniem dla wydajności systemu okazało się generowanie opisów przez lokalny model LLM. Początkowa wersja kodu zakładała przetwarzanie wierszy pojedynczo: program tworzył wiersz, wysyłał zapytanie do serwera Ollama, czekał na odpowiedź i dopiero wtedy przechodził do kolejnego rekordu. Rozwiązanie to było bardzo powolne.

\subsection*{Identyfikacja problemu}
Pomiary wykazały, że wygenerowanie 100 recenzji produktów przez model Llama 3 zajmowało blisko 130 sekund. Prawie 95\% tego czasu stanowiło oczekiwanie na odpowiedź sztucznej inteligencji, podczas gdy zwykłe dane tekstowe (np. nazwiska z biblioteki Faker) generowały się w ułamki sekund. Czas oczekiwania na jeden rekord wynosił około 1.3 sekundy, a procesor oraz karta graficzna przez większość tego czasu nie były w pełni obciążone.

\subsection*{Rozwiązanie dwuetapowe}
W celu lepszego wykorzystania zasobów obliczeniowych, wprowadzono algorytm dwufazowy:

\begin{enumerate}
    \item \textbf{Faza 1 -- dane proste}: W pierwszej kolejności aplikacja generuje wszystkie standardowe wiersze dla tabeli, całkowicie pomijając kolumny LLM. Działa to natychmiastowo i tworzy gotowe obiekty, które mogą posłużyć jako kontekst dla zapytań AI.
    \item \textbf{Faza 2 -- grupowe zapytania LLM}: Zebrane konteksty są następnie przekazywane w paczkach do oddzielnej funkcji, która zajmuje się wyłącznie komunikacją z modelem.
\end{enumerate}

\subsection*{Grupowanie zapytań}
Kluczem do przyspieszenia procesu okazało się grupowanie zapytań w mniejsze pakiety. Zamiast wysyłać pojedynczy wiersz, program opiera się na parametrze \texttt{SUB\_BATCH\_SIZE} (domyślnie 15) i wysyła do modelu kilkanaście kontekstów w jednym dużym prompcie. W odpowiedzi model zwraca tekstową strukturę \gls{JSON} zawierającą od razu cały pakiet wyników. Znacznie skraca to narzut czasu potrzebny na wielokrotne inicjalizowanie zapytań przez silnik AI.

Dodatkowo, zastosowano mechanizm \texttt{asyncio.Semaphore}, który zezwala serwerowi Ollama na przetwarzanie kilku takich paczek jednocześnie (domyślnie 8 zapytań w tle). Jeśli wystąpi błąd formatowania dla danej paczki (np. model pominie znak w odpowiedzi \gls{JSON}), to ten mały fragment uszkodzonych 15 wierszy jest wyłapywany i przerabiany bezpiecznym, pojedynczym zapytaniem jeden po drugim, by nie psuć całego postępu generowania.

\subsection*{Testy wydajności}
W celu dobrania optymalnych wielkości paczek przeprowadzono testy wydajnościowe. Jako stację wykorzystano laptopa z kartą graficzną NVIDIA RTX 5060, korzystając z modelu Llama 3 uruchomionego lokalnie w środowisku Ollama dla próbki 100 wierszy AI.

\begin{table}[H]
    \centering
    \begin{tabular}{|c|c|c|c|}
        \hline
        \textbf{Numer próby} & \textbf{Czas [s]} & \textbf{Czas na 1 wiersz [s]} & \textbf{Zmiana w konfiguracji}       \\
        \hline
        1                    & 129,52            & 1,30                          & Generowanie pojedyncze (sekwencyjne) \\
        \hline
        2                    & 218,88            & 2,19                          & Zbyt wysoka współbieżność zapytań    \\
        \hline
        3                    & 63,43             & 0,63                          & Implementacja pakietowania (batches) \\
        \hline
        4                    & 124,41            & 1,24                          & Testowanie innych rozmiarów limitów  \\
        \hline
        5                    & 52,24             & 0,52                          & Optymalnie dopasowane parametry      \\
        \hline
    \end{tabular}
    \caption{Czasy generowania 100 wierszy AI na laptopie RTX 5060}
    \label{tab:llm_optimization_100}
\end{table}

Zgodnie z Tabelą \ref{tab:llm_optimization_100}, przy najlepiej dobranych parametrach udało się skrócić czas ze 130 sekund do około 52. Oznacza to blisko 2.5-krotne przyspieszenie skryptu bez negatywnego wpływu na złożoność całego programu. Zaobserwowano również, że zbyt wysokie ustawienie równoległych żądań (Próba 2) doprowadziło do przeciążenia API i drastycznego wydłużenia pracy.

\subsection*{Alternatywny silnik -- vLLM}

Początkowe prace nad aplikacją oraz testy lokalne opierały się na serwerze Ollama \cite{ollama} -- narzędziu zaprojektowanym z myślą o prostocie obsługi i łatwym uruchamianiu modeli LLM na komputerach osobistych. Ollama wyróżnia się intuicyjnym interfejsem i prostym procesem instalacji, co doskonale sprawdza się przy deweloperskich eksperymentach. Jednakże w zastosowaniach chmurowych, wymagających dużej przepustowości i serwowania wielu równoczesnych zapytań, jego architektura nie oferuje optymalnej wydajności.

W celu maksymalizacji prędkości generowania w środowisku chmurowym zdecydowano się przeprowadzić testy z wykorzystaniem vLLM \cite{vllm} -- biblioteki serwującej modele LLM, opracowanej na Uniwersytecie Kalifornijskim w Berkeley. vLLM różni się od Ollama podejściem architektonicznym: zamiast interfejsu zorientowanego na użytkownika końcowego, oferuje zaawansowany silnik zoptymalizowany pod kątem produkcyjnej wydajności i maksymalnego wykorzystania sprzętu.

Kluczowym mechanizmem vLLM jest \textit{PagedAttention} -- algorytm zarządzania pamięcią GPU inspirowany mechanizmami stronicowania pamięci z systemów operacyjnych. Standardowe podejście do obsługi mechanizmu uwagi \cite{attention_ml} w modelach LLM wymaga alokacji ciągłego bloku pamięci dla każdego zapytania, co prowadzi do fragmentacji i marnotrawstwa zasobów GPU. PagedAttention dzieli pamięć kontekstu KV-cache (pamięć podręczna klucz-wartość) na małe, nieciągłe bloki, które mogą być dynamicznie przydzielane i zwalniane. Dzięki temu vLLM osiąga blisko optymalną utylizację pamięci karty graficznej i pozwala na obsługę znacznie większej liczby równoległych żądań.

Drugą istotną zaletą jest mechanizm \textit{continuous batching}, który dynamicznie łączy napływające zapytania w czasie rzeczywistym, bez konieczności czekania na skompletowanie całej paczki. Pozwala to na pełniejsze wykorzystanie mocy obliczeniowej kart GPU, gdyż silnik nieustannie przetwarza kolejne tokeny z różnych zapytań zamiast pozostawać bezczynny między kolejnymi żądaniami.

\subsection*{Konfiguracja vLLM na platformie Kaggle}

Wdrożenie vLLM w środowisku Kaggle wymagało rozwiązania dodatkowych problemów środowiskowych. Domyślna wersja vLLM (0.17) korzysta z biblioteki FlashInfer, która wymaga kompilacji kerneli CUDA w czasie wykonywania programu. System plików Kaggle, będący w dużej części środowiskiem tylko do odczytu, uniemożliwiał tę kompilację, powodując błędy budowania oprogramowania.

Rozwiązaniem okazało się zastosowanie starszej wersji vLLM (0.6.6), która zamiast FlashInfer wykorzystuje bibliotekę xFormers \cite{xformers}. xFormers dostarcza prekompilowane, wydajne implementacje operacji uwagi, co eliminuje konieczność kompilacji kerneli CUDA w czasie uruchamiania serwera.

Model Llama 3 8B-Instruct uruchomiono z podziałem tensorów na dwie karty Tesla T4 (\texttt{tensor-parallel-size=2}), co pozwoliło na wykorzystanie łącznie 32 GB pamięci VRAM. Aplikacja, dzięki oparciu komunikacji na standardzie OpenAI API, wymagała jedynie zmiany adresu URL serwera modelu oraz wyłączenia sztucznego podziału na paczki -- dotychczasowy interfejs i cała logika przetwarzania po stronie backendu pozostały niezmienione, delegując pełne zarządzanie ruchem do silnika vLLM.

\subsection*{Skalowalność dla dużych baz danych}

W celu ostatecznego porównania wydajności wykonano próbę obciążeniową polegającą na wygenerowaniu dużej bazy danych zawierającej aż 10\,000 wierszy przetworzonych przez LLM. Testy przeprowadzono na dwóch platformach sprzętowych oraz z dwoma silnikami inferencji (Ollama i vLLM).

\begin{table}[H]
    \centering
    \begin{tabular}{|l|l|c|c|c|}
        \hline
        \textbf{Środowisko testowe} & \textbf{Silnik} & \textbf{Wierszy AI} & \textbf{Czas [s]} & \textbf{Czas/wiersz [s]} \\
        \hline
        Laptop (RTX 5060)           & Ollama          & 100                 & 52,24             & 0,52                     \\
        \hline
        Laptop (RTX 5060)           & Ollama          & 10\,000             & 8\,038,12         & 0,80                     \\
        \hline
        Kaggle (2x Tesla T4)        & Ollama          & 100                 & 178,02            & 1,78                     \\
        \hline
        Kaggle (2x Tesla T4)        & Ollama          & 10\,000             & 14\,860,23        & 1,49                     \\
        \hline
        Kaggle (2x Tesla T4)        & vLLM            & 100                 & 18,30             & 0,18                     \\
        \hline
        Kaggle (2x Tesla T4)        & vLLM            & 10\,000             & 1\,693,50         & 0,17                     \\
        \hline
    \end{tabular}
    \caption{Porównanie skalowalności w zależności od sprzętu i silnika}
    \label{tab:llm_scalability}
\end{table}

Wygenerowanie 10\,000 wierszy na laptopie z serwerem Ollama zajęło ponad 8\,038 sekund (nieco ponad 2 godziny). Długotrwałe obciążenie karty graficznej wiązało się ze zjawiskiem tzw. \textit{thermal throttlingu} (zabezpieczenie przed przegrzaniem obniżające taktowanie procesora), przez co średni czas utworzenia jednego wiersza wzrósł z 0,52 do 0,80 sekundy. Mimo to aplikacja zniosła ten test poprawnie i zapisała całą tabelę w pliku bazy SQLite.

Dla porównania, ten sam schemat z serwerem Ollama uruchomiono w darmowym środowisku chmurowym Kaggle, dysponującym dwoma kartami Tesla T4 (łącznie 32 GB pamięci VRAM). Pełna próba 10\,000 wierszy na Kaggle zajęła aż 14\,860 sekund (ponad 4 godziny). Różnica wydajności względem laptopa wynika ze specyfikacji sprzętowych -- karty Tesla T4 oparte są na architekturze Turing \cite{turing_microarchitecture} z 2018 roku, która nie posiada ulepszonych jednostek obliczeniowych obecnych w nowoczesnych kartach z serii RTX (architektura Ada Lovelace \cite{ada_lovelace_microarchitecture}).

Zastąpienie serwera Ollama przez vLLM na tej samej platformie Kaggle przyniosło drastyczną poprawę wydajności. Generowanie 100 wierszy AI skróciło się ze 178 do zaledwie 18,3 sekundy (przyspieszenie blisko 10-krotne), a pełna próba 10\,000 wierszy ze 14\,860 do 1\,693 sekund (około 28 minut zamiast ponad 4 godzin) -- przyspieszenie blisko 9-krotne. Co więcej, średni czas na wiersz z vLLM (0,17 s) okazał się ponad 4-krotnie lepszy niż najlepsza lokalna konfiguracja Ollama na nowszym laptopie RTX 5060 (0,52s). Wynik ten potwierdza, że mechanizmy PagedAttention i \textit{continuous batching} zaimplementowane w vLLM pozwalają na znacznie efektywniejsze wykorzystanie chmurowych zasobów GPU niż tradycyjne podejście serwera Ollama.

\section{Odporność na błędy łączności z API}
Komunikacja z zewnętrznymi interfejsami, takimi jak API usługi OpenAI, niesie ze sobą ryzyko sporadycznej utraty połączenia sieciowego lub przeciążeń po stronie dostawcy chmury. W kodzie serwera uodporniono wywołania API za pomocą precyzyjnych bloków wychwytywania wyjątków \texttt{try-except}. W przypadku wystąpienia \textit{timeoutu} (braku odpowiedzi sieciowej dla konkretnego zapytania), system nie odrzuca całego procesu po wielu godzinach pracy. Kod zachowuje odpowiedni tekst błędu (np. \texttt{Connection error}) wpisując go bezpiecznie w odpowiednią komórkę bazy danych tabeli. Pozłwala to na płynne kontynuowanie generowania kolejnego wiersza uniemożliwiając restart skryptu, pozostawiając operatorowi wgląd do omyłkowego elementu tuż poprzez odpowiednie odnalezienie słowa błędu w ostatecznym arkuszu docelowym Excel bądź CSV.

